\section*{الفصل \ref{ch:g7:earthquakes} --- الزلازل}

\begin{solution}{exo:g7:earthquakes:1}
(أ) \omterm{def:g7:earthquakes:focus}{البؤرة}؛ (ب) \omterm{def:g7:earthquakes:intensity}{الشدة}؛ (ج) \omterm{def:g7:earthquakes:fault}{الصدع}؛ (د) \omterm{def:g7:earthquakes:focus}{المركز السطحي}؛ (ه) \omterm{def:g7:earthquakes:magnitude}{القدر}؛ (و)
\omterm{def:g7:earthquakes:seismometer}{مقياس الزلازل}.
\end{solution}

\begin{solution}{exo:g7:earthquakes:2}
(أ) عند القرية نفسها: \omterm{def:g7:earthquakes:focus}{فالمركز السطحي} هو نقطة السطح الواقعة فوق \omterm{def:g7:earthquakes:focus}{البؤرة}
مباشرة. (ب) \qty{30}{km}، رأسيًا إلى الأسفل. (ج) أدنى عادة: فالاهتزاز يضعف
مع المسافة. (وهو يتوقف أيضًا على الأرض: فعلى \omterm{def:g3:stones-pebbles-sand:sand}{الوحل} الطري، قد تهتز القرية
الثانية مع ذلك بقوة.)
\end{solution}

\begin{solution}{exo:g7:earthquakes:3}
(أ) \omterm{def:g7:earthquakes:magnitude}{قدر} (عدد برقم عشري واحد). (ب) \omterm{def:g7:earthquakes:intensity}{شدة} (ما رُئي في مكان واحد). (ج) \omterm{def:g7:earthquakes:intensity}{شدة}
(أرقام رومانية، مختلفة في مكانين). (د) \omterm{def:g7:earthquakes:magnitude}{قدر}: الطاقة المتحررة، أكبر بدرجة
واحدة من \omterm{def:g7:earthquakes:magnitude}{القدر}.
\end{solution}

\begin{solution}{exo:g7:earthquakes:4}
(أ) خطأ: العلماء لا يستطيعون التنبؤ \omterm{def:g4:volcanoes-earthquakes-first-look:earthquake}{بالزلازل}؛ وكل ما يعرفونه أي الصدوع
ستنكسر من جديد يومًا ما. (ب) خطأ: في عرض \omterm{def:g2:water-on-earth:ocean}{المحيط} نادرًا ما يزيد علو \omterm{def:g7:earthquakes:tsunami}{التسونامي}
على متر واحد؛ وهو يعلو قرب \omterm{def:g1:landscapes-around-us:water}{الساحل}. (ج) صواب. (د) صواب: 92~\% من \omterm{def:g4:volcanoes-earthquakes-first-look:earthquake}{الزلازل} التي
\omterm{def:g7:earthquakes:magnitude}{قدرها}~6 أو أكثر تقع على بعد أقل من \qty{300}{km} من حد بين \omterm{def:g7:plate-tectonics:plate}{الصفائح}.
\end{solution}

\begin{solution}{exo:g7:earthquakes:5}
(أ) نحو 127 في السنة \omterm{def:g7:earthquakes:magnitude}{قدرها} من 6 إلى 6.9؛ ونحو 13 \omterm{def:g7:earthquakes:magnitude}{قدرها} من 7 إلى 7.9.
(ب) $126.7 \div 13.12 \approx 9.7$: أقل بنحو 10~مرات. (ج)
$13.12 + 0.76 + 0.04 = 13.92$: نحو \textbf{14} في السنة. (د)
$0.04 \times 50 = 2$: اثنان في 50~سنة (سومطرة 2004 واليابان 2011).
\end{solution}

\begin{solution}{exo:g7:earthquakes:6}
(أ) درجة واحدة: 32~مرة. (ب) درجتان: $32 \times 32 = 1024$، أي نحو
1000~مرة. (ج) كل درجة تجعل الأمواج أكبر 10~مرات: $10 \times 10 =
100$~مرة.
\end{solution}

\begin{solution}{exo:g7:earthquakes:7}
(أ) $8 \times 5 = \qty{40}{km}$. (ب) $8 \times 25 = \qty{200}{km}$. (ج)
$120 \div 8 = \qty{15}{s}$. (د) المحطة الواحدة لا تعطي إلا مسافة، لا
اتجاهًا: فكل النقاط الواقعة على تلك المسافة تشكل دائرة حول المحطة، وقد
يكون \omterm{def:g7:earthquakes:focus}{المركز السطحي} في أي مكان عليها.
\end{solution}

\begin{solution}{exo:g7:earthquakes:8}
(أ) V. (ب) II. (ج) VIII. (د) \omterm{def:g7:earthquakes:intensity}{الشدة} تتغير من مكان إلى آخر: البلاغ (ج)
يأتي من مكان قريب من \omterm{def:g7:earthquakes:focus}{المركز السطحي}، أو مبني على أرض طرية؛ والبلاغ (أ) من
مكان أبعد، أو على \omterm{def:g3:stones-pebbles-sand:rock}{صخر} صلب. ومع ذلك يبقى \omterm{def:g4:volcanoes-earthquakes-first-look:earthquake}{للزلزال} \omterm{def:g7:earthquakes:magnitude}{قدر} واحد فقط.
\end{solution}

\begin{solution}{exo:g7:earthquakes:9}
(أ) $7031 - 5561 = 1470$. (ب) $1470 \div 7031 \approx 0.209$: نحو
\textbf{21~\%}. (ج) \omterm{def:g7:plate-tectonics:plate}{الصفيحة} الغائصة تنزل عميقًا في \omterm{def:g6:inside-the-earth:layers}{الوشاح}. وهي لا تزال
صلبة بما يكفي لتنكسر، فيمكن أن تبدأ \omterm{def:g4:volcanoes-earthquakes-first-look:earthquake}{الزلازل} داخلها، على عمق مئات
الكيلومترات.
\end{solution}

\begin{solution}{exo:g7:earthquakes:10}
(أ) $2500 \div 800 = 3.125$~ساعة: 3~ساعات و$0.125 \times 60 = 7.5$
دقيقة، أي نحو 3~ساعات و8~دقائق. (ب) 3~ساعات و8~دقائق $-$ ساعتان
و17~دقيقة $=$ 51~دقيقة من التأخر. (ج) استغرقت الموجة ساعات لتصل، بينما
يُحدَّد مكان \omterm{def:g4:volcanoes-earthquakes-first-look:earthquake}{زلزال} كبير تحت \omterm{def:g2:water-on-earth:ocean}{البحر} في غضون دقائق: فكان يمكن للإنذار أن يرسل
الناس إلى الأراضي المرتفعة قبل مجيء الموجة بوقت طويل.
\end{solution}

\begin{solution}{exo:g7:earthquakes:11}
(أ) $400 \div 800 = 0.5$~ساعة: 30~دقيقة. (ب) $20 \div 40 = 0.5$~ساعة:
30~دقيقة. (ج) 60~دقيقة. (د) آخر \qty{20}{km} تستغرق ما تستغرقه أول
\qty{400}{km}: فالموجة تتباطأ تباطؤًا هائلًا في المياه الضحلة. (ه) مقدمة
الموجة تتباطأ بينما الماء الذي خلفها لا يزال يأتي مسرعًا: فيتراكم الماء،
وتزداد الموجة علوًا.
\end{solution}

\begin{solution}{exo:g7:earthquakes:12}
(أ) ثلاث درجات: $32 \times 32 \times 32 = 32\,768$، أي نحو
\textbf{33\,000}~مرة. (ب) $32\,768 \div 127 \approx 258$: نحو
260~سنة من كل \omterm{def:g4:volcanoes-earthquakes-first-look:earthquake}{زلازل} العالم التي \omterm{def:g7:earthquakes:magnitude}{قدرها} 6. (ج) تكاد الطاقة كلها تتحرر في
القليل من أكبر \omterm{def:g4:volcanoes-earthquakes-first-look:earthquake}{الزلازل} على الإطلاق: \omterm{def:g4:volcanoes-earthquakes-first-look:earthquake}{فزلزال} واحد \omterm{def:g7:earthquakes:magnitude}{قدره}~9 يفوق قرونًا من
\omterm{def:g4:volcanoes-earthquakes-first-look:earthquake}{الزلازل} التي \omterm{def:g7:earthquakes:magnitude}{قدرها} 6.
\end{solution}

\begin{solution}{pb:g7:earthquakes:1}
\textbf{1.} A: $45 - 25 = 20$~ثانية. B: $29 - 16 = 13$~ثانية. C: $34 - 19 =
15$~ثانية.

\textbf{2.} المحطة B: تأخرها هو الأقصر، وكلما قصر التأخر كان \omterm{def:g4:volcanoes-earthquakes-first-look:earthquake}{الزلزال}
أقرب. (ووصلت أمواجها الأسرع أولًا كذلك.)

\textbf{3.} A: $8 \times 20 = \qty{160}{km}$. B: $8 \times 13 =
\qty{104}{km}$. C: $8 \times 15 = \qty{120}{km}$.

\textbf{4.} كل النقاط الواقعة على بعد \qty{160}{km} من A: دائرة نصف قطرها
\qty{160}{km} مركزها A. \omterm{def:g7:earthquakes:focus}{والمركز السطحي} في مكان ما عليها.

\textbf{5.} A: $160 \div 20 = 8$~مربعات. B: $104 \div 20 = 5.2$~مربع.
C: $120 \div 20 = 6$~مربعات.

\textbf{6.} تقع E على بعد $3 + 6 = 9$~مربعات شرقًا و14.5~مربعًا شمالًا من
O، في عرض \omterm{def:g2:water-on-earth:ocean}{البحر} بين المحطات الثلاث؛ ودائرتا A وB تمران بها.

\textbf{7.} تحت \omterm{def:g2:water-on-earth:ocean}{البحر}. فالحركة المفاجئة لقاع \omterm{def:g2:water-on-earth:ocean}{البحر} قد تطلق \omterm{def:g7:earthquakes:tsunami}{تسونامي}: ولا بد
من إنذار \omterm{def:g1:landscapes-around-us:water}{السواحل} والجزيرة.

\textbf{8.} درجة واحدة من \omterm{def:g7:earthquakes:magnitude}{القدر}: طاقة أكبر بنحو 32~مرة.

\textbf{9.} $12 \times 20 = \qty{240}{km}$.

\textbf{10.} $240 \div 800 = 0.3$~ساعة، و$0.3 \times 60 = 18$~دقيقة.

\textbf{11.} بعد 5~دقائق من \omterm{def:g4:volcanoes-earthquakes-first-look:earthquake}{الزلزال}.

\textbf{12.} $18 - 5 = 13$: أمام أهل بورت~هيل نحو \textbf{13~دقيقة}
ليبلغوا أرضًا مرتفعة. (تتباطأ الموجة في المياه الضحلة قرب الجزيرة، وهذا
يعطي وقتًا أطول قليلًا، لكن لا ينبغي لأحد أن يعتمد على ذلك: فهم يغادرون
في الحال.)
\end{solution}
